\subsection{Formal depth and descent of the torsion class}
\label{subsec:conductor-picard-descent}

The exact order of the logarithmic form also measures the order to
which the torsion line bundle is trivial on the infinitesimal
neighborhoods of the curve. At conductor equality, this gives a
precise obstruction to lifting the radial action to the line bundle.
A $p$th tensor power supplies a lift; at an external center
its descent again decreases the conductor.

\begin{proposition}\label{prop:torsion-formal-depth}
Let $X$ be a normal projective integral surface over an algebraically
closed field $k$ of characteristic $p>0$, let $Y\subset X$ be a
smooth integral effective Cartier divisor contained in the smooth
locus, with $Y^2>0$, and let $M\ne\OO_X$ have order $p$
with $M|_Y\simeq\OO_Y$.
Let $pa$ be the exact order along $Y$ of its logarithmic form.
Then
\[
 M|_{aY}\simeq\OO_{aY},
 \qquad M|_{(a+1)Y}\not\simeq\OO_{(a+1)Y}.
\]
In particular, $a$ is the largest integer $n$ for which $M|_{nY}$
is trivial.
\end{proposition}

\begin{proof}
Put $I=\OO_X(-Y)$. For every $n\geq1$ one has
\begin{equation}\label{eq:thickening-constant-functions}
 H^0(nY,\OO_{nY})=k.
\end{equation}
Indeed, the successive kernels in the restriction to smaller
thickenings are
$I^j/I^{j+1}\simeq\OO_Y(-jY)$ for $j\geq1$.
Their degrees are negative, so they have no nonzero sections.
Starting with $H^0(Y,\OO_Y)=k$ proves the assertion; constants
provide the surviving sections. Consequently
\begin{equation}\label{eq:thickening-cohomology-injection}
 H^1(X,I^n)\longrightarrow H^1(X,\OO_X)
 \quad\text{is injective}
\end{equation}
for every $n\geq1$.

We first establish a local approximation. Choose a sufficiently
small smooth affine neighborhood of a point of $Y$ with
\'etale coordinates $t,z$, where $Y=(t=0)$, and denote its ring
by $R$. The monomials $t^iz^j$, $0\leq i,j<p$, form a basis of
$R$ over $R^p$. This is the polynomial-ring calculation transported
along the \'etale map by relative Frobenius
\cite[Tag~0EBS]{Stacks}. If $u\in R^*$ and
$du\in t^{pa}\Omega^1_{R/k}$, write
\[
 u=\sum_{0\leq i,j<p}r_{ij}^{\,p}t^iz^j.
\]
In the $dt$ coefficient of $du$, the monomials with $i>0$ are
distinct members of the $p$-basis; the $dz$ coefficient treats
those with $j>0$. Since $t^{pa}\in R^p$, freeness implies that
$r_{ij}^{\,p}$ is divisible by $t^{pa}$ whenever $(i,j)\ne(0,0)$.
Thus $r_{ij}$ is divisible by $t^a$, and
\begin{equation}\label{eq:local-torsion-pth-approximation}
 u=v^p+t^{pa}b.
\end{equation}
Here $v=r_{00}$ is a unit after shrinking, because $u$ is a unit
along $Y$. This calculation takes place in the algebraic local
rings.

Choose local frames $e_i$ of $M$ and a trivialization of $M^p$,
and let $u_i$ be the corresponding units. Then
$\omega=du_i/u_i$. Replacing the frames using
\eqref{eq:local-torsion-pth-approximation} gives, near $Y$,
\[
 u_i=1+t_i^{pa}b_i.
\]
If $e_i=g_{ij}e_j$, then
$g_{ij}^p=u_i/u_j\equiv1\pmod{I^{pa}}$.
Hence $(g_{ij}-1)^p\in I^{pa}$, and regularity of the parameter
defining $Y$ gives $g_{ij}-1\in I^a$.
The frames therefore glue on $aY$, proving the first assertion.

Suppose that $M$ were trivial on $(a+1)Y$. Lift a trivializing
frame locally, so that $g_{ij}\equiv1\pmod{I^{a+1}}$.
The units $u_i$ then agree modulo $I^{p(a+1)}$, and glue to a
global unit on $p(a+1)Y$. By
\eqref{eq:thickening-constant-functions} this unit is a scalar.
Rescaling the frames by a $p$th root of that scalar gives
\[
 u_i=1+t_i^{p(a+1)}b_i.
\]
Differentiation produces no term of order $p(a+1)-1$, so
$\omega$ has order at least $p(a+1)$ along $Y$. This contradicts
its exact order $pa$.
\end{proof}

\begin{proposition}\label{prop:radial-linearization-obstruction}
Assume the hypotheses and equality of
Theorem~\ref{thm:conductor-equality-structure}.
Let $\delta$ be the radial derivation induced by
$x_0\partial/\partial x_0$ in the homogeneous coordinates used there, and
write $I=\OO_X(-Y)$.
For transition functions $g_{ij}$ of $M$, put
\[
 o_\delta(M)=
 \bigl[\delta(g_{ij})/g_{ij}\bigr]\in H^1(X,\OO_X).
\]
Then $M$ admits a $\mu_p$-linearization for the radial action if
and only if $o_\delta(M)=0$.
Moreover,
\begin{equation}\label{eq:radial-leading-obstruction}
 p\nmid a\quad\Longrightarrow\quad o_\delta(M)\ne0.
\end{equation}
Thus a linearization, and in particular descent of $M$ to the
radial quotient, requires $p\mid a$.
\end{proposition}

\begin{proof}
A lift of $\delta$ to $M$ is a $k$-linear operator $\nabla$ with
$\nabla(fe)=\delta(f)e+f\nabla(e)$.
In the frame $e_i$, write $\nabla(e_i)=\alpha_i e_i$.
Compatibility is precisely
\[
 \alpha_i-\alpha_j=\delta(g_{ij})/g_{ij}.
\]
Thus the displayed class is independent of the frames and is the
obstruction to an ordinary lift.

If a lift exists, $\nabla^p-\nabla$ is $\OO_X$-linear:
its commutator with multiplication by $f$ is multiplication by
$\delta^p(f)-\delta(f)=0$.
It is therefore multiplication by a scalar $c\in k$.
Choose $\lambda\in k$ such that $\lambda^p-\lambda=-c$.
The adjusted operator satisfies
\[
 (\nabla+\lambda)^p=\nabla+\lambda.
\]
The eigenspace projections for the distinct roots of $T^p-T$
are polynomials in this operator. They commute with restriction
to open sets, and the Leibniz rule makes the resulting character
grading compatible with the grading on $\OO_X$ induced by
$\delta$. This is a $\mu_p$-linearization.
Conversely, differentiating a linearization gives such a lift,
with the restricted identity already satisfied.
This proves the equivalence.

By Proposition~\ref{prop:torsion-formal-depth}, the first
nontrivial infinitesimal class of $M$ is
\[
 0\ne\xi\in H^1(Y,I^a/I^{a+1}).
\]
Indeed, the square-zero unit sequence for
$aY\subset(a+1)Y$ identifies the kernel of restriction on Picard
groups with this cohomology group: the global units on both
thickenings are $k^*$ and restriction is surjective.
Choose frames compatible with the trivialization on $aY$.
Their transitions satisfy $g_{ij}-1\in I^a$ and represent $\xi$
modulo $I^{a+1}$.

The derivation vanishes on $Y$ and has weight one on $I/I^2$.
It therefore has weight $a$ on $I^a/I^{a+1}$, so
\[
 \frac{\delta(g_{ij})}{g_{ij}}
 \equiv a(g_{ij}-1)\pmod{I^{a+1}}.
\]
The obstruction is represented by a cocycle in $I^a$, whose image
in $H^1(Y,I^a/I^{a+1})$ is $a\xi$.
When $p\nmid a$, it is nonzero. Its image in $H^1(X,\OO_X)$
remains nonzero by
\eqref{eq:thickening-cohomology-injection}.
This proves \eqref{eq:radial-leading-obstruction}.
A bundle pulled back from the quotient has its canonical
linearization, giving the final assertion.
\end{proof}



\begin{proposition}\label{prop:external-conductor-descent}
Assume the hypotheses and equality of
Theorem~\ref{thm:conductor-equality-structure}, and suppose that
its radial center $v$ lies outside $X$. Let
$q:X\to S=X/\mu_p$ be the radial quotient and let $Y_S=q(Y)$.
Then $S$ is smooth projective, $q|_Y:Y\to Y_S$ is an isomorphism,
and $Y_S$ is a smooth very ample divisor satisfying
\begin{equation}\label{eq:external-quotient-polarization}
 q^*Y_S=pY,\qquad Y_S^2=pY^2.
\end{equation}
The homomorphism $q^*:\Pic(S)\to\Pic(X)$ is injective.

If $M$ admits a $\mu_p$-linearization, it descends to a nontrivial
line bundle $M_S$ of exact order $p$, trivial on $Y_S$.
The exact order of its logarithmic form along $Y_S$ is $a$.
Thus $p\mid a$ and
\[
 2g(Y_S)-2=p(a/p)Y_S^2,\qquad a/p\geq5.
\]
\end{proposition}

\begin{proof}
Smoothness and projectivity of $S$ follow from the quotient
calculation in
Theorem~\ref{thm:conductor-equality-structure}.
To identify its polarization, let $R$ be the homogeneous
coordinate ring of $X$, with $\mu_p$ scaling $x_0$ and fixing
$x_1,\ldots,x_r$. The invariant ring is a quotient of
\[
 k[z,x_1,\ldots,x_r],\qquad z=x_0^p,\quad
 \deg z=p,\quad\deg x_i=1.
\]
Indeed, an equivariant surjection is surjective on each character
summand. The $p$th Veronese of this weighted polynomial ring
is generated in degree one: a monomial of weight $np$ is a
product of $n$ monomials of weight $p$. The same holds for its
quotient. Hence $S=\operatorname{Proj}(R^{\mu_p})$ is embedded
by $\OO_S(p)$.

There is an isomorphism
\[
 R^{\mu_p}/(x_0^p)\simeq R/(x_0).
\]
Surjectivity follows from the invariant generators $x_i$, $i>0$.
An invariant element divisible by $x_0$ is $x_0$ times an element
of weight $p-1$. Lift the latter to a polynomial of that weight;
each of its monomials is divisible by $x_0^{p-1}$ with invariant
quotient. This proves the kernel assertion.
Thus $z$ cuts out $Y_S$ scheme-theoretically, identifies it with
$Y$, and gives $\OO_S(Y_S)=\OO_S(p)$.
Its pullback is the divisor of $x_0^p$, proving the first
identity in \eqref{eq:external-quotient-polarization}.
The second follows by intersection and $\deg q=p$.
This argument does not require projective normality of $X$.

We next prove injectivity on Picard groups. If $q^*P$ is trivial,
its canonical linearization acts on its one-dimensional space
of global sections by a character. A nonzero section is nowhere
zero. Evaluate it at a $k$-point of the fixed divisor $Y$.
The action on that fibre of a pullback bundle is trivial, so the
character is trivial. The invariant section descends, since
\[
 (q_*q^*P)^{\mu_p}=P
\]
by the projection formula and
$(q_*\OO_X)^{\mu_p}=\OO_S$.
Surjectivity of $q$ makes the descended section nowhere zero.
Therefore $P$ is trivial.

Suppose now that $M$ is linearized. Its action on the trivial
bundle $M|_Y$ is one character, since $Y$ is proper integral.
Twist the linearization by the inverse character so that its
fibre action on $Y$ is trivial.
On an invariant affine neighborhood of a point of $Y$, a
weight-zero section of $M$ lifts a nonzero vector in the fibre:
take the weight-zero part of any lift.
After shrinking, it is an invariant frame.
Away from $Y$, the function $x_0/x_i$ is an invertible function
of weight one on suitable charts. A homogeneous local frame
exists by decomposing a frame into its character components
and choosing one nonzero at the point. Multiplication by a
power of $x_0/x_i$ makes its weight zero.
The nonvanishing opens are invariant and descend to opens of
$S$, since $q$ is a universal homeomorphism.
Thus invariant frames glue to a line bundle
\[
 M_S=(q_*M)^{\mu_p},\qquad q^*M_S\simeq M.
\]
Injectivity of $q^*$ makes $M_S$ nontrivial of exact order $p$,
and $Y\simeq Y_S$ gives $M_S|_{Y_S}\simeq\OO_{Y_S}$.

Let $\omega_S$ be its logarithmic form. Its $p$th-power
trivialization pulls back to one differing from that of $M$
by a scalar, so
\begin{equation}\label{eq:descended-logarithmic-form}
 q^*\omega_S=\omega.
\end{equation}
In particular $\omega_S\ne0$.
Let $n$ be its exact order along $Y_S$. The restriction to
$\Omega^1_{Y_S}$ is zero, since $M_S|_{Y_S}$ is trivial.
The remaining conormal section has negative degree, so $n\geq1$.
The leading tangential coefficient is nonzero: otherwise the
leading section would lie in
\[
 H^0\bigl(Y_S,\OO_{Y_S}(-(n+1)Y_S)\bigr)=0.
\]
Write locally $\omega_S=s^n(A\,ds+B\,dz)$ with $B$ nonzero
modulo $s$. Closedness has leading coefficient
$n s^{n-1}B\,ds\wedge dz$, so $p\mid n$.

At a general closed point of $Y$, the completed quotient
coordinates are $s=t^p$, $z=z$. Pullback kills $ds$, and the
nonzero tangential term has exact order $pn$. Thus
\[
 pa=pn,\qquad n=a.
\]
The use of completed local rings at a closed point makes no
perfectness assumption on $k(Y)$.
Consequently the new conductor integer is $a/p$, and
\eqref{eq:external-quotient-polarization} gives
\[
 2g(Y_S)-2=pa\,Y^2=aY_S^2=p(a/p)Y_S^2.
\]
Theorem~\ref{thm:conductor-equality-structure}, applied to the new
smooth datum, gives $a/p\geq5$.
\end{proof}

\begin{lemma}\label{lem:radial-power-descent}
Under the hypotheses and equality of
Theorem~\ref{thm:conductor-equality-structure}, suppose
$M\simeq L^{\otimes p}$ for an invertible sheaf $L$ on $X$.
Then $M$ admits a radial $\mu_p$-linearization, $p\mid a$,
and $M=q^*P$ for an invertible sheaf $P$ on the radial quotient
$q:X\to S$. If the center is external, $P$ is the nontrivial
order-$p$ descended bundle of
Proposition~\ref{prop:external-conductor-descent}.
\end{lemma}

\begin{proof}
Let $g_{ij}$ be the transition functions of $L$.
The transitions $g_{ij}^p$ of $M$ are invariant under the
$\mu_p$-coaction. Indeed, writing its coordinate algebra as
$k[\epsilon]/(\epsilon^p)$, the difference between the coaction
of a function and that function lies in the ideal $(\epsilon)$,
whose elements have $p$th power zero.
Every open is stable under the height-one action.
Giving the local tensor-power frames weight zero therefore
glues a linearization of $M$.
Proposition~\ref{prop:radial-linearization-obstruction} gives
$p\mid a$.

For the invertible descent, on every invariant affine chart
one has $R^p\subset R^{\mu_p}\subset R$.
Since $k$ is perfect, relative Frobenius factors as
\[
 X\xrightarrow{q}S\xrightarrow{g}X^{(p)},
 \qquad gq=F_{X/k}.
\]
The ring inclusions glue; moreover $R^{\mu_p}$ is finite over
$R^p$ as a submodule of the finite $R^p$-module $R$.
Set $P=g^*(L^{(p)})$. Then
\[
 q^*P=F_{X/k}^*L^{(p)}\simeq L^{\otimes p}\simeq M.
\]
For an external center, injectivity of $q^*$ in
Proposition~\ref{prop:external-conductor-descent} identifies
$P$ with that proposition's descended bundle.
\end{proof}

The radial descent uses a very ample embedding, whereas the
normalization argument in Section~\ref{subsec:positive-square-picard}
uses only positive self-intersection. The finite-root obstruction
therefore applies also to ample curves and to curves on normal
surfaces that are smooth in their neighborhood. For a general
N\'eron--Severi torsion class the necessary roots may not exist;
the obstruction in Corollary~\ref{thm:positive-square-picard}
states exactly which divisibility is excluded at equality.
